\subsection{Train versus prediction}

Una vez congelada la tabla, el primer análisis no supervisado comparó la distribución de las
138 filas etiquetadas con las 59 filas objetivo. Se examinaron 1,401 variables numéricas. Para
cada feature se calcularon, entre otros, diferencia estandarizada de medias, missingness y
fracción del conjunto objetivo fuera del rango observado de entrenamiento.

Una forma genérica del standardized mean difference es
\begin{equation}
\mathrm{SMD}_j
=
\frac{\bar x_{j,\Train}-\bar x_{j,\Target}}
{\sqrt{\frac{s^2_{j,\Train}+s^2_{j,\Target}}{2}}}.
\label{eq:smd}
\end{equation}
Sólo seis variables superaron al menos un threshold descriptivo. Cuatro correspondían a
resúmenes históricos Sentinel-2 de agosto, una a la desviación de densidad aparente
SoilGrids 30--60 cm y una a EVI mínimo Sentinel-2 2025.

\begin{table}[H]
\centering
\caption{Principales flags descriptivos del diagnóstico train--prediction.}
\label{tab:cp02-shift}
\small
\begin{tabular}{p{0.60\linewidth}rr}
\toprule
Feature & $|\mathrm{SMD}|$ & fuera de rango \\
\midrule
Sentinel-2 ENDDI histórico, agosto & 0.5343 & 0.0000 \\
Sentinel-2 DSWI2 histórico, agosto & 0.5303 & 0.0000 \\
Sentinel-2 MCRC histórico, agosto & 0.5296 & 0.0169 \\
Sentinel-2 CRC histórico, agosto & 0.5035 & 0.0339 \\
SoilGrids bdod 30--60 cm, std & 0.3218 & 0.1017 \\
Sentinel-2 EVI mínimo 2025, std & 0.2212 & 0.1017 \\
\bottomrule
\end{tabular}
\end{table}

El resultado sugería ausencia de un shift univariado masivo, pero no descartaba shift
multivariado. Esa pregunta fue retomada formalmente con clasificación adversarial en
Checkpoint~04A.

\subsection{Ablaciones estructuradas}

En lugar de entrenar inmediatamente sobre las 1,401 variables numéricas, Checkpoint~02 definió
una secuencia acumulativa de ocho representaciones:

\begin{table}[H]
\centering
\caption{Ablaciones de Checkpoint 02.}
\label{tab:cp02-ablations}
\small
\begin{tabular}{llr}
\toprule
ID & Contenido & $p$ \\
\midrule
A0 & geometría + administración & 12 \\
A1 & A0 + ambiente estático & 115 \\
A2 & A1 + historia SIAP & 129 \\
A3 & A2 + clima histórico & 153 \\
A4 & A3 + BASIC histórico & 692 \\
A5 & clean completo & 692 \\
A6 & clean + remote/weather 2025, sin SIAP outcome & 1,396 \\
A7 & competition completo, incluido SIAP 2025 & 1,401 \\
\bottomrule
\end{tabular}
\end{table}

A4 y A5 coinciden numéricamente porque, bajo el manifest de esa versión, todas las features
clean numéricas ya estaban contenidas en A4. La duplicación se mantuvo como checkpoint
semántico: A4 representa hasta BASIC histórico y A5 todo clean.

\subsection{Por qué las ablaciones importan en high-dimensional small-\texorpdfstring{$n$}{n}}

Cuando $p\gg n$, agregar una familia completa puede degradar un modelo incluso si contiene
señal real. La varianza del ajuste, colinealidad y ruido pueden dominar. Por ello se mide
\begin{equation}
\Delta\RMSE_{A\rightarrow B}
=
\RMSE(B)-\RMSE(A)
\label{eq:ablation-delta}
\end{equation}
por modelo y protocolo, en vez de suponer monotonicidad con el número de features.

El ejemplo más claro aparece con Ridge: en state-stratified, A3 obtiene RMSE medio 0.5298,
mientras A4/A5 suben a 0.6707. Agregar cientos de resúmenes satelitales a un modelo lineal de
alpha fijo no es una mejora automática. ExtraTrees es mucho menos sensible en ese protocolo:
A3 da 0.5035 y A7 0.4966.

\subsection{Clean y competition como ejes experimentales}

Checkpoint~02 mostró que la señal 2025 completa podía ayudar en algunos modelos, pero también
que SIAP contemporáneo no tenía por qué mejorar de forma automática. ExtraTrees pasa de
0.5001 en A6 a 0.4966 en A7 bajo state-stratified, una ganancia pequeña; en
municipality-grouped, A6=0.8835 y A7=0.8916, es decir, el agregado completo empeora. Este
patrón anticipa la conclusión mucho más fuerte de Checkpoint~04E.1: una fuente contemporánea
puede ser legal y plausible sin ser un buen proxy parcelario.

\subsection{Rango y extrapolación}

La fracción de objetivos fuera del rango train,
\[
q_j=
\frac{1}{|\Target|}
\sum_{i\in\Target}
\mathbf 1\{x_{ij}<\min_{\Train}x_{\cdot j}\;\vee\;
x_{ij}>\max_{\Train}x_{\cdot j}\},
\]
es útil como alerta, pero no prueba que una parcela sea imposible de predecir. Un punto puede
estar fuera del rango de una feature marginal y, simultáneamente, cerca de entrenamiento en
una representación multivariada relevante. Por esa razón el soporte final de Checkpoint~04
combina geografía, embeddings, temporalidad y clasificación adversarial.

\begin{technicalnote}
Los flags de covariate shift de Checkpoint~02 son descriptivos. No se usaron para inferir
pseudo-etiquetas ni para eliminar variables usando conocimiento de las 59 respuestas.
\end{technicalnote}
